\begin{exercise}{33}
Let $X$ be a simplicial set and consider the canonical map $X\to\mathrm N(hX)$.
\begin{enumerate}[label=(\arabic*),leftmargin=*,itemsep=0pt,parsep=0pt,
                  topsep=.5\baselineskip]
 \item Show that this map factors through the canonical map
       $X\to\mathrm{cosk}_2(X)$.
 \item Show that the induced map $\mathrm{cosk}_2(X)\to\mathrm N(hX)$
       is an isomorphism if $X$ is isomorphic to the nerve of a category.
 \item Show that the map $\mathrm{cosk}_2(X)\to\mathrm N(hX)$ is in general
       not an isomorphism. Hint: Find an $X$ which is $2$-coskeletal, but not
       the nerve of a category.
 \item Prove or disprove the following statement: The map
       $\mathrm{cosk}_2(X)\to\mathrm N(hX)$ is an isomorphism if and only if
       $X$ is isomorphic to the nerve of a category.
\end{enumerate}

\begin{proof}
Write $C=\mathrm{cosk}_2$ and $u_Z\colon Z\to CZ$ for the
canonical map. Let $\eta_X\colon X\to\mathrm N(hX)$ be the
unit of Proposition~1.2.18.

(1) By Exercise~\ref{ex:19}, $u_{\mathrm N(hX)}$ is an
isomorphism. Define the natural comparison
\[
 v_X=(u_{\mathrm N(hX)})^{-1}C(\eta_X)
 \colon CX\longrightarrow\mathrm N(hX).
\]
Naturality of $u$ gives $v_Xu_X=\eta_X$, as required.

(2) If $X=\mathrm N(\mathcal D)$, then $u_X$ is an isomorphism
by Exercise~\ref{ex:19}. Also $\eta_X$ is an isomorphism:
the counit $h\mathrm N(\mathcal D)\to\mathcal D$ is the
isomorphism of Corollary~1.2.14, and the adjunction triangle
identity identifies $\eta_{\mathrm N(\mathcal D)}$ with its
inverse after applying $\mathrm N$.
Thus $v_X=\eta_Xu_X^{-1}$ is an isomorphism. These properties
transport across an isomorphism $X\cong\mathrm N(\mathcal D)$.

(3) Take $X=I^2$. By Definition~1.1.38, $CX$ agrees with
$X$ through degree $2$, so it has no edge from $0$ to $2$.
By Exercise~\ref{ex:26}, $hX=[2]$, hence
$\mathrm N(hX)=\Delta^2$. The comparison is the identity on
vertices, and therefore misses the edge $0\to2$.
For the hint, $CI^2$ is itself $2$-coskeletal by
Exercise~\ref{ex:20}. It has edges $0\to1$ and $1\to2$ but
no edge $0\to2$, so it cannot be a nerve.

(4) Part (2) proves one direction. The converse fails for
$X=\partial\Delta^3$.
The inclusion $j\colon X\hookrightarrow\Delta^3$ is an
isomorphism through degree $2$. Thus $Cj$ is an isomorphism
by Definition~1.1.38, and $hj$ is an isomorphism by
Lemma~1.2.7. Naturality of $v$ gives
\[
 \mathrm N(hj)v_X=v_{\Delta^3}Cj.
\]
Since $\Delta^3=\mathrm N([3])$, part (2) makes
$v_{\Delta^3}$ an isomorphism, so the displayed equality
proves that $v_X$ is an isomorphism.
However, $\partial\Delta^3$ is not a nerve by
Exercise~\ref{ex:16}.
\end{proof}
\end{exercise}
